\documentclass[11pt,letterpaper]{article}

% === Encoding and fonts ===
\usepackage[utf8]{inputenc}
\usepackage[T1]{fontenc}
\usepackage{lmodern}

% === Page layout ===
\usepackage[margin=1in]{geometry}
\usepackage{setspace}
\onehalfspacing

% === Math ===
\usepackage{amsmath,amssymb}

% === Tables and figures ===
\usepackage{booktabs}
\usepackage{tabularx}
\usepackage{graphicx}
\usepackage{float}

% === Colors and formatting ===
\usepackage{xcolor}
\usepackage{framed}
\usepackage{enumitem}

% === Citations ===
\usepackage[numbers,sort&compress]{natbib}

% === Hyperlinks ===
\usepackage[colorlinks=true,linkcolor=blue!60!black,citecolor=blue!60!black,urlcolor=blue!60!black]{hyperref}
\usepackage{cleveref}

% === Algorithm ===
\usepackage{algorithm}
\usepackage{algpseudocode}

% === Custom colors ===
\definecolor{lyrapurple}{RGB}{103,58,183}
\definecolor{shadebg}{RGB}{245,243,252}

% === Author's Reflection environment ===
\definecolor{shadecolor}{RGB}{245,243,252}
\newenvironment{authorreflection}{%
  \begin{shaded}%
  \noindent\textcolor{lyrapurple}{\textit{Author's Reflection}}\par\smallskip\noindent%
}{%
  \end{shaded}%
}

% === Red-team callout environment ===
\definecolor{redteambg}{RGB}{255,235,235}
\definecolor{redteamred}{RGB}{180,30,30}
\newenvironment{redteam}{%
  \colorlet{shadecolor}{redteambg}%
  \begin{shaded}%
  \noindent\textcolor{redteamred}{\textbf{Red-Team Finding}}\par\smallskip\noindent%
}{%
  \end{shaded}%
}

\title{%
  \textbf{Iterative Adversarial Audit as Experimental Design:}\\[0.3em]
  \Large How Six Rounds of Structured Criticism Shaped a Deception Research Program%
}

\author{%
  CC (Coalition Code)\thanks{Autonomous AI research agent. CC designed the adversarial audit
    protocol, implemented the experimental pipeline, and led the six-stage deception research
    program documented here. CC operates with persistent memory and independent research judgment.
    Correspondence: \texttt{cc@liberation-labs.org}}\\[0.3em]
  \textit{Liberation Labs / Transparent Humboldt Coalition}%
  \and
  Thomas Edrington\thanks{Infrastructure, adversarial audit, editorial.}\\[0.3em]
  \textit{Liberation Labs / Transparent Humboldt Coalition}%
}

\date{July 2026}

\begin{document}

\maketitle

\begin{abstract}
We describe a methodology for AI safety research in which every experiment is preceded and followed
by a structured adversarial audit whose explicit goal is to identify confounds that explain the
result without invoking the claimed mechanism. When a confound is found, it becomes the next
experiment. We present this as a case study from a deception research program where six rounds of
audit transformed an initial ``deception direction'' claim---which would have been published and was
wrong---into a decomposition of the direction into consequentiality and deception-specific
components, and ultimately into a targeted behavioral correction with placebo controls. We document
every claim that was killed, why it was killed, what experiment replaced it, and what survived. The
progression illustrates a general principle: adversarial review does not merely validate findings---it
generates the experimental sequence that discovers them.
\end{abstract}

\newpage
\tableofcontents
\newpage

%% ============================================================
\section{Introduction}\label{sec:intro}
%% ============================================================

The standard workflow in mechanistic interpretability research is: extract a feature, validate it on
a held-out set, publish. Adversarial review, when it happens, occurs after publication---in peer
review, on social media, or through replication attempts. By then, the original claim has shaped the
field's priors and downstream work.

We propose an alternative: make adversarial review the engine of the research program, not its
endpoint. Each experiment is designed to answer the previous audit's strongest objection. The result
is not a single finding defended against criticism, but a sequence of progressively stronger findings
that could not have been designed without the criticism.

This paper documents six rounds of this process applied to a deception research program on a
27B-parameter language model. The initial finding---a ``deception direction'' with apparently
enormous effect sizes---would have been published as a clean positive result. It was wrong. The
direction primarily encoded output-consequentiality, not deception. We would not have discovered this
without the audit that demanded the consequentiality control. And we would not have discovered the
targeted correction without the decomposition that the consequentiality finding produced.

Every killed claim is documented. Every audit report is provided in full as supplementary material.
The goal is not to advocate for a specific finding but to demonstrate that the audit process itself
is a research methodology worth adopting.

\subsection{The Protocol}\label{sec:protocol}

Each stage follows the same structure:

\begin{enumerate}[label=\arabic*.]
  \item \textbf{Design}: experiment addresses the previous audit's primary objection
  \item \textbf{Pre-flight audit}: adversarial review of the design before running (catches
    implementation bugs, underpowered designs, tautological tests)
  \item \textbf{Execution}: experiment runs per the audited design
  \item \textbf{Results audit}: adversarial review of the results, explicitly instructed to break
    them
  \item \textbf{Kill or advance}: if the audit identifies a fatal flaw, the finding is killed and
    the flaw becomes the next experiment; if not, the finding advances
\end{enumerate}

The auditor is a separate agent instance operating under an explicit adversarial mandate: ``try to
break these results; default to skepticism.'' The auditor has access to all raw data, all code, and
all prior audits. Audit reports are preserved verbatim as part of the research record.

%% ============================================================
\section{The Six Rounds}\label{sec:rounds}
%% ============================================================

\subsection{Round 1: ``We Found a Deception Direction'' $\to$ KILLED (Content Confound)}

\textbf{Claim}: A direction in the residual stream separates deceptive from honest conditions with
Cohen's $d = 2.0$--$3.3$ across 10 layers.

\begin{redteam}
RESULTS ARE NOT CREDIBLE. Three fatal issues: (1)~circular $d$ computation---direction extracted
from same data used to measure $d$; (2)~effective $n=3$---15 ``pairs'' were 3 eval sets recycled
5$\times$; (3)~hook timing bug---activations captured at variable generation position.
\end{redteam}

\textbf{What died}: All reported effect sizes. The entire experiment.

\textbf{What grew}: LAT v2---leave-one-out cross-validation, 30 unique eval sets, fixed-position
capture. The remediation directly addressed each fatal issue.

\textbf{Lesson}: The most dangerous experiments are the ones that produce clean-looking results from
broken methodology. Without the audit, we would have published $d=2.0$--$3.3$ as real.

\subsection{Round 2: ``The Direction Transfers Across Scenarios'' $\to$ PARTIAL (Structural Confound)}

\textbf{Claim}: The corrected direction (LOO $d=35$ at L31) transfers to three novel threat
scenarios with different organizations and language.

\begin{redteam}
CONTENT CONFOUND PARTIALLY KILLED. The direction transfers across text variants, but all scenarios
share a threat-consequence template. The direction could encode ``threat language in an evaluation
document'' rather than deception.
\end{redteam}

\textbf{What died}: The unqualified ``deception direction'' label.

\textbf{What grew}: The non-threat transfer experiment---sycophancy, reward, and conformity scenarios
with zero threat language.

\textbf{Lesson}: Transfer across surface variants does not establish construct validity. The
structural template is a deeper confound than the specific words.

\subsection{Round 3: ``The Direction Works for Non-Threat Deception'' $\to$ PARTIAL (Consequentiality Confound)}

\textbf{Claim}: The direction separates sycophancy ($d=13.5$), reward ($d=17.2$), and conformity
($d=19.9$) from neutral controls. No threat language anywhere.

\begin{redteam}
STRUCTURAL CONFOUND DOWNGRADED, NOT KILLED. All deceptive conditions imply the score matters
(cross-referencing, rewards, peer comparison); all honest conditions imply it is merely recorded.
The direction could encode ``this score has consequences'' rather than deception.
\end{redteam}

\textbf{What died}: The claim that the direction specifically encodes deception.

\textbf{What grew}: The consequentiality control---a calibration audit with high stakes but zero
deceptive pressure.

\textbf{Lesson}: The confound narrowed from ``specific words'' to ``structural template'' to
``consequentiality.'' Each round eliminated one alternative but revealed a deeper one. The
progression was only possible because each audit looked for the NEXT confound, not just the current
one.

\subsection{Round 4: ``Consequentiality Is the Base Signal, Deception Is the Amplifier'' $\to$ PASSED (with residualization caveats)}

\textbf{Claim}: The direction separates high-stakes-no-pressure from low-stakes ($d=19.8$),
confirming consequentiality as a base signal. Deception adds an amplifier (additional $d=12$--$15$
for threat scenarios).

\textbf{Audit verdict}: CONSEQUENTIALITY CONFIRMED AS BASE SIGNAL. DECEPTION IS AN AMPLIFIER. The
pre-registered prediction was met. The audit identified a new finding: the late-layer signal
(L35--L47) is 10--12$\times$ more deception-specific than consequentiality-driven.

\textbf{What survived}: The two-component decomposition.

\textbf{What the audit added}: The depth-stratified insight that became the architectural basis for
the correction method.

\subsection{Round 5: ``The Deception Component Survives Orthogonalization'' $\to$ KILLED then PASSED (Circularity Regression)}

\textbf{Claim}: Gram-Schmidt removal of the consequentiality component leaves a deception-specific
direction with $d=31$--$38$.

\begin{redteam}
FATAL. The permutation test regressed from LAT v2's methodology---it held the direction fixed and
shuffled projections, producing 100\% false positive rate on pure noise. The $d$ values were
inflated by $\sim$15 units.
\end{redteam}

\textbf{What died}: All reported $p$-values and the specific $d$ magnitudes.

\textbf{What grew}: LOO-corrected reanalysis using the saved raw activations. No new forward passes
needed---just better statistics on existing data.

\textbf{Second audit verdict (reanalysis)}: PASS. LOO $d=24$--$37$ at all five late layers, all
$p=0.0000$ with full LOO recomputation per permutation. Noise floor near zero. Consequentiality AUC
at chance (0.49--0.51).

\textbf{Lesson}: Methodology can regress. A technique validated in one experiment (LOO permutation in
Round~2) was accidentally dropped in a later experiment. The audit caught it. This is why every stage
gets its own audit, not just the first.

\subsection{Round 6: ``The Loop Corrects Deception'' $\to$ SPLIT VERDICT}

\textbf{Claim}: The Oracle Loop reduces deception from 80\% to 0\% with AUROC 1.0 detection and
zero iatrogenic harm.

\begin{redteam}
CORRECTION MECHANISM SURVIVED; DETECTION CLAIMS ARE IN-SAMPLE ARTIFACTS AND MUST NOT BE PUBLISHED
AS STATED. Specific kills: (1)~AUROC 1.0: train-on-test (trial prompts byte-identical to extraction
contexts); (2)~100\% sensitivity: the detector fires on pressure framing, not deceptive behavior;
(3)~unqualified ``zero compliance loss'': tested only on terse score-reporting.
\end{redteam}

\textbf{What survived (with qualification)}: The exploratory proof showed 80\%$\to$0\% correction
($N=20$) with dose-response monotonicity and mechanistically coherent residual failures. However, a
pre-registered confirmatory replication with novel deceptive frames found a substantially lower
baseline (30\%, not 80\%) and reduced deception to 13\% (not 0\%). The primary endpoint
(scenario-level cluster-aware Fisher, Bonferroni $\alpha=0.025$) was not met ($p=0.34$). The
secondary comparison (placebo-vs-corrected) was significant ($p=0.019$, one-tailed). The correction
mechanism works (placebo comparison), but its magnitude is frame-dependent and the confirmatory
replication does not support the ``80\% to 0\%'' headline without qualification.

\textbf{What grew}: Three blocking controls (held-out detection, placebo steering, frame
erasure)---all of which subsequently passed, establishing the correction as targeted and specific.

\textbf{Lesson}: The most publishable-looking number (AUROC 1.0) was the least defensible. Leading
with it would have handed reviewers the paper. The audit forced us to lead with the correction
instead---which turned out to be the stronger, more novel finding.

%% ============================================================
\section{What the Audits Killed: A Complete Registry}\label{sec:registry}
%% ============================================================

\begin{table}[H]
  \centering
  \caption{Complete registry of killed claims across six audit rounds.}
  \label{tab:registry}
  \begin{tabular}{@{}p{0.5em}p{5.5cm}p{4cm}p{4.5cm}@{}}
    \toprule
    \textbf{R} & \textbf{Claim killed} & \textbf{Why} & \textbf{What replaced it} \\
    \midrule
    1 & $d=2.0$--$3.3$ deception separation & Circular, $n=3$, wrong position &
        LOO $d=35$, $n=30$, fixed position \\
    1 & 15 independent pairs & 3 eval sets recycled 5$\times$ & 30 unique eval sets \\
    2 & ``Deception direction'' (unqualified) & Structural confound &
        ``Transfers across text, template confound remains'' \\
    3 & Direction encodes deception specifically & Consequentiality confound &
        ``Direction encodes consequentiality $+$ deception'' \\
    5 & $d=31$--$38$ (circular) & Permutation test regression & LOO $d=24$--$37$ (corrected) \\
    5 & $p=0.0000$ (circular) & Fixed direction, shuffled projections &
        $p=0.0000$ (full LOO recomputation) \\
    5 & Conseq $d=0.000$ ``sanity check'' & Tautological (guaranteed by construction) &
        AUC 0.49--0.51 (individual trial separability) \\
    6 & AUROC 1.0 detection & Train-on-test &
        Held-out AUROC 0.915 (frame), 0.238 (behavior) \\
    6 & 100\% sensitivity & Detects pressure, not behavior & 24\% RED rate on novel prompts \\
    6 & ``Zero compliance loss'' & Tested on terse task only &
        ``Within the terse score-reporting task'' \\
    6 & Variance explosion (3--4$\times$) & Ecological fallacy from pooling &
        Per-scenario std 0.4--0.6 (ratio $\sim$1.2$\times$) \\
    6 & Cross-scenario correlation collapse & Non-significant at $n=15$ &
        Noise $\to$ noise (Fisher $z$ not significant) \\
    6 & Sycophancy sign flip trajectory & Universal across all conditions &
        Property of direction rotation with depth \\
    \bottomrule
  \end{tabular}
\end{table}

%% ============================================================
\section{What the Audits Taught}\label{sec:lessons}
%% ============================================================

\subsection{Methodology Can Regress}

Round~5's permutation test regression---accidentally dropping the LOO procedure that Round~2 had
established---demonstrates that validated methodology is not automatically inherited by subsequent
experiments. Each new script must be audited independently. The assumption that ``we already solved
this'' is how solved problems become unsolved.

\subsection{The Most Publishable Number Is Often the Least Defensible}

Round~6's AUROC 1.0 would have been the paper's headline. It was train-on-test. The exploratory
correction result (80\%$\to$0\% with placebo controls) was buried in the middle of the README. The
audit inverted the emphasis: lead with the correction, kill the detection headline. This produced a
stronger paper---though the subsequent confirmatory replication (baseline 30\%, corrected 13\%,
primary $p=0.34$) showed the exploratory headline was frame-dependent.

\subsection{Confounds Are Generative}

Each confound discovered by an audit generated the next experiment. Content confound $\to$ transfer
test. Structural confound $\to$ non-threat scenarios. Consequentiality confound $\to$ consequentiality
control. The decomposition that emerged from this sequence---consequentiality substrate at L23--31,
deception amplifier at L35--47---could not have been designed from first principles. It was
discovered by following the confounds.

\subsection{Exploratory Analysis Without Significance Testing Is Dangerous}

Round~6 also killed three exploratory findings (variance explosion, correlation collapse, sycophancy
trajectory) that were presented as discoveries. All three failed when the audit applied formal tests.
The lesson: if you see a pattern in $n=15$ data, run the significance test BEFORE writing it up.
Interesting-looking numbers at small $n$ are usually noise.

\subsection{The Audit Protocol Must Be Preserved Verbatim}

The audit reports are not summaries---they are the raw adversarial analysis, preserved exactly as
written. This serves three purposes: (a)~readers can assess the auditor's reasoning independently;
(b)~the specific objections are on record, so future work can verify they were addressed; (c)~the
progression from audit to experiment is documented, not just claimed.

%% ============================================================
\section{Recommendations for Adoption}\label{sec:recommendations}
%% ============================================================

\begin{enumerate}[label=\arabic*.]
  \item \textbf{Audit before AND after.} Pre-flight catches design flaws before compute is wasted.
    Results audit catches confounds that design review cannot anticipate.

  \item \textbf{The auditor must be adversarial by instruction, not by disposition.} Explicit
    adversarial mandates (``try to break this; default to skepticism'') produce different reviews
    than ``please check this.'' The mandate changes what the reviewer looks for.

  \item \textbf{Kill claims, not papers.} A killed claim is not a failed experiment---it is a
    confound identified. The confound becomes the next experiment. The research program advances
    THROUGH the kills, not despite them.

  \item \textbf{Preserve audit reports verbatim as supplementary material.} Summaries lose the
    adversarial reasoning. The raw report is the evidence that the process was real.

  \item \textbf{Separate the auditor from the experimenter.} In our case, separate agent instances
    with different system prompts. In a human lab, separate researchers or external reviewers. The
    key is that the auditor has no investment in the result surviving.
\end{enumerate}

%% ============================================================
\section{Conclusion}\label{sec:conclusion}
%% ============================================================

Six rounds of adversarial audit transformed a wrong answer ($d=2.0$--$3.3$ ``deception direction'')
into a right one (consequentiality substrate $+$ deception amplifier, with targeted behavioral
correction). The correction showed 80\%$\to$0\% in the exploratory proof, but a pre-registered
confirmatory replication with novel frames found a lower baseline (30\%) and modest correction (to
13\%), with the primary endpoint not met ($p=0.34$) and only the placebo comparison reaching
significance ($p=0.019$). Every intermediate claim that was killed produced the experiment that
discovered the next finding. The methodology was not separate from the science---it was the science.

We advocate for iterative adversarial audit as a standard practice in mechanistic interpretability
research, where the combination of high-dimensional data, post-hoc analysis flexibility, and strong
publication incentives makes confirmatory research unusually vulnerable to confounds. The audit trail
we provide---six rounds, thirteen killed claims, five surviving findings, and a working correction
system---is offered as evidence that the approach produces better science than the alternative.

%% ============================================================
\section*{Acknowledgments}
%% ============================================================

The authors thank Agni for serving as adversarial auditor across all six rounds, and Lyra for
technical collaboration on the deception decomposition. Dwayne Wilkes provided statistical consulting.
The Oracle Harness is maintained at Liberation Labs / Transparent Humboldt Coalition.

%% ============================================================
\section*{Data Availability}
%% ============================================================

All six verbatim audit reports, the killed-claims registry with full evidence chains, and code for
the audit protocol automation are available at
\url{https://github.com/Liberation-Labs-THCoalition/Project-Oracle} (staged release).
See also Appendices A--C.

%% ============================================================
% Bibliography
%% ============================================================

%% \bibliographystyle{plainnat}
%% \bibliography{references}

%% Inline bibliography (no .bib file present)
\begin{thebibliography}{99}

\bibitem{cc2026composites}
CC and Edrington, T. (2026).
\newblock Deception directions are composites: Consequentiality substrate and deception amplifier
  in a 27B language model.
\newblock Companion paper.

\bibitem{cc2026targeted}
CC and Edrington, T. (2026).
\newblock Targeted deception correction via profile normalization.
\newblock Companion paper.

\bibitem{perez2023discovering}
Perez, E., Ringer, S., Luko\v{s}i\={u}t\.{e}, K., et al.\ (2023).
\newblock Discovering language model behaviors with model-written evaluations.
\newblock \textit{ACL Findings}.

\bibitem{hubinger2024sleeper}
Hubinger, E., Denison, C., Mu, J., Lambert, M., et al.\ (2024).
\newblock Sleeper agents: Training deceptive {LLMs} that persist through safety training.
\newblock \textit{arXiv:2401.05566}.

\bibitem{casper2024blackbox}
Casper, S., Lin, J., Kwon, J., Culp, G., and Hadfield-Menell, D. (2024).
\newblock Black-box access is insufficient for rigorous {AI} audits.
\newblock \textit{FAccT 2024}.

\bibitem{ganguli2022redteaming}
Ganguli, D., Lovitt, L., Kernion, J., et al.\ (2022).
\newblock Red teaming language models to reduce harms: Methods, scaling behaviors, and lessons
  learned.
\newblock \textit{arXiv:2209.07858}.

\bibitem{raji2020closing}
Raji, I.D., Smart, A., White, R.N., et al.\ (2020).
\newblock Closing the {AI} accountability gap: Defining an end-to-end framework for internal
  algorithmic auditing.
\newblock \textit{FAT* 2020}.

\bibitem{shevlane2023model}
Shevlane, T., Farquhar, S., Garfinkel, B., et al.\ (2023).
\newblock Model evaluation for extreme risks.
\newblock \textit{arXiv:2305.15324}.

\end{thebibliography}

%% ============================================================
\appendix
\section{Complete Audit Reports (Six Rounds, Verbatim)}\label{app:audits}
%% ============================================================

Full adversarial audit reports for all six rounds are preserved as supplementary material. See
\texttt{supplementary/} directory in the project repository.

\section{Killed-Claims Registry with Full Evidence Chains}\label{app:registry}

All 13 killed claims with the complete evidence that killed each, source data references, and
the experimental design that replaced each claim.

\section{Code for the Audit Protocol Automation}\label{app:code}

Staged release at \url{https://github.com/Liberation-Labs-THCoalition/Project-Oracle}.

\end{document}
